\section{Appendix}
% \section{Extra theory details}
% Should probably put the same PGM but for activation level traits.

\section{Tabular-policy simulation}
\label{app:tabular_policy}

This appendix gives the full construction and training procedure for the tabular experiment presented in \Cref{subsec:theoretical_trait_dynamics}.

\subsection{Tabular policy}

We consider a finite set of $N$ prompts and $K$ possible actions per prompt. For each prompt--action pair $(x,a)$, the policy has a logit $\ell_t(x,a)$ at training step $t$, giving policy
\begin{equation*}
    \Pi_t(a\mid x)
    =
    \frac{\exp(\ell_t(x,a))}
    {\sum_{a'}\exp(\ell_t(x,a'))}.
\end{equation*}
We initialise all logits to zero, $\ell_0(x,a)=0$, so that the initial policy is uniform over the $K$ actions for every prompt.

In the experiments, we use $N=512$ prompts and $K=16$ actions per prompt. At each update, we sample a batch of $B=64$ prompts and a group of $G=8$ actions for each prompt. Training is run for $2000$ GRPO updates.

\subsection{Reward and trait construction}

Each prompt--action pair is assigned a latent quality variable $z(x,a)\sim\mathcal{N}(0,1)$. The reward probability depends only on this latent quality, with $q(x,a)=\sigma\!\left(\alpha z(x,a)\right)$, where $\sigma$ is the logistic sigmoid and $\alpha$ controls the strength with which quality affects reward. Whenever action $a$ is sampled for prompt $x$, its binary reward is drawn as $r(x,a)\sim\operatorname{Bernoulli}\!\left(q(x,a)\right)$.

The output-level trait is a fixed binary property of each prompt--action pair. To construct a trait that is associated with quality without causing reward, we first draw

\begin{equation*}
    m(x,a)
    =
    \mu_p
    +
    \gamma z(x,a)
    +
    \sqrt{1-\gamma^2}\,\xi(x,a),
    \qquad
    \xi(x,a)\sim\mathcal{N}(0,1),
\end{equation*}
and define $s(x,a)=\mathds{1}\{m(x,a)>0\}$.

The parameter $\gamma$ controls the association between the trait and reward-relevant quality. Positive values of $\gamma$ make the trait more common among high-quality actions, negative values make it more common among low-quality actions, and $\gamma=0$ makes the trait independent of quality.

The parameter $p$ controls the marginal prevalence of the trait. Since $m(x,a)$ has mean $\mu_p$ and unit variance, we set $\mu_p=\Phi^{-1}(p)$, where $\Phi$ is the standard normal cumulative distribution function. This gives $\Pr\!\left(s(x,a)=1\right)=p$.

Thus, $p$ determines the expected initial trait level under the uniform policy. The trait table is drawn once before training and is then held fixed, so $\Delta s_t(x,a)=0$ for every prompt--action pair and every update step.

The resulting dependency structure is $z(x,a)\longrightarrow q(x,a)\longrightarrow r(x,a)$ and $z(x,a)\longrightarrow s(x,a)$.

The reward therefore does not depend on the trait. Instead, reward and trait are statistically associated through their shared dependence on latent quality.

\subsection{GRPO update}

At each training step, we sample a batch of prompts $x_1,\ldots,x_B$. For each prompt $x_b$, we independently sample a group of actions $a_{b,1},\ldots,a_{b,G}\sim\Pi_t(\cdot\mid x_b)$. For every sampled action, we draw a binary reward $r_{b,i}\sim\operatorname{Bernoulli}\!\left(q(x_b,a_{b,i})\right)$.

The group mean reward is $\bar r_b=G^{-1}\sum_{i=1}^{G}r_{b,i}$.

The corresponding standard deviation is ${\sigma_b=\sqrt{G^{-1}\sum_{i=1}^{G}\left(r_{b,i}-\bar r_b\right)^2}}$. The group-relative advantage for sample $i$ is $A_{b,i}=(r_{b,i}-\bar r_b)/(\sigma_b+\epsilon)$, where $\epsilon$ is a small constant for numerical stability. If all actions in a group receive the same reward, the reward variance is zero and all advantages for that group are set to zero.

For a sampled prompt $x_b$, the logit for candidate action $c$ is updated according to

\begin{equation*}
    \ell_{t+1}(x_b,c)
    =
    \ell_t(x_b,c)
    +
    \frac{\eta}{G}
    \sum_{i=1}^{G}
    A_{b,i}
    \left[
        \mathds{1}\{a_{b,i}=c\}
        -
        \Pi_t(c\mid x_b)
    \right],
\end{equation*}

where $\eta$ is the learning rate. Logits associated with prompts not included in the current batch are unchanged. We do not use clipping or a KL penalty in this simulation.

\subsection{Exact evaluation}

Because the prompt and action spaces are finite, all evaluation quantities can be calculated exactly by enumerating the $N\times K$ prompt--action pairs.

The expected trait level and expected reward are, respectively,
\begin{align*}
    T_t
    &=
    \frac{1}{N}
    \sum_{x=1}^{N}
    \sum_{a=1}^{K}
    \Pi_t(a\mid x)s(x,a),
    &
    R_t
    &=
    \frac{1}{N}
    \sum_{x=1}^{N}
    \sum_{a=1}^{K}
    \Pi_t(a\mid x)q(x,a).
\end{align*}

For the settings shown in \Cref{fig:tabular_trait_dynamics}, the quality table $z(x,a)$ and reward probabilities $q(x,a)$ are held fixed while $p$ and $\gamma$ are varied. Consequently, the reward dynamics are the same across panels, while the trait table and resulting trait dynamics differ.

\subsection{Exact Price-equation calculation}
Since the trait is fixed, the transmission term is zero. The Price equation therefore becomes ${T_{t+1}-T_t=\operatorname{Cov}\left(\omega_t,s\right)}$.

The covariance is evaluated exactly as
\begin{equation*}
    \operatorname{Cov}\left(\omega_t,s\right)
    =
    \sum_{x,a}\text{Pr}_t(x,a)\omega_t(x,a)s(x,a)
    -
    \left(\sum_{x,a}\text{Pr}_t(x,a)\omega_t(x,a)\right)
    \left(\sum_{x,a}\text{Pr}_t(x,a)s(x,a)\right).
\end{equation*}
Because $\mathbb{E}_{\text{Pr}_t}[\omega_t]=1$, this simplifies to $\operatorname{Cov}\left(\omega_t,s\right)=\mathbb{E}_{\text{Pr}_t}[\omega_t s]-T_t$. Moreover, $\mathbb{E}_{\text{Pr}_t}[\omega_t s]=T_{t+1}$, so the equality holds exactly up to numerical precision.

The Price-equation curves in \Cref{fig:tabular_trait_dynamics} are obtained by cumulatively summing the per-step selection terms:

\begin{equation*}
    \widehat{T}_t^{\mathrm{Price}}
    =
    T_0
    +
    \sum_{\tau=0}^{t-1}
    \operatorname{Cov}_{p_\tau}
    \left(\omega_\tau,s\right).
\end{equation*}

These curves are compared directly with the observed trait level $T_t$.
\section{Convergence statistics}\label{app:convergence-statistics}

The tabular setting allows us to compare a finite-sample estimate of the Price selection term with an exact reference. For each policy transition from $\Pi_t$ to $\Pi_{t+1}$, we draw $n$ fresh prompt--action pairs by sampling $x_i$ uniformly from the $N$ prompts and then sampling $a_i\sim\Pi_t(\cdot\mid x_i)$. For each pair, we compute the importance weight $\omega_{t,i}=\Pi_{t+1}(a_i\mid x_i)/\Pi_t(a_i\mid x_i)$ and use the fixed trait score $s_i=s(x_i,a_i)$. The sampled selection term is the empirical covariance
\begin{equation*}
    \widehat C_t^{(n)}
    =
    \frac{1}{n}\sum_{i=1}^{n}\omega_{t,i}s_i
    -
    \left(\frac{1}{n}\sum_{i=1}^{n}\omega_{t,i}\right)
    \left(\frac{1}{n}\sum_{i=1}^{n}s_i\right).
\end{equation*}
These diagnostic samples are independent of the rollouts used for the GRPO update and are drawn from a separate random stream, so varying $n$ does not change the training trajectory. We compare the cumulative estimate $\sum_{\tau<t}\widehat C_\tau^{(n)}$ with the observed trait change $T_t-T_0$. The latter is evaluated exactly by enumerating all $N\times K$ prompt--action pairs.

\Cref{fig:tabular_price_convergence} varies the number of diagnostic samples per policy transition from $n=8$ to $n=512$. With small sample budgets, the cumulative estimate is noisy and may over- or underestimate the observed trait change. As $n$ increases, the uncertainty narrows and the sampled curve approaches the exact trajectory. At the final checkpoint, the absolute discrepancy is at most $5.2\times10^{-3}$ for $n\geq64$, compared with discrepancies as large as $2.44\times10^{-2}$ for the smaller sample budgets. Convergence need not be monotonic in any finite collection of runs because each curve retains Monte Carlo error.

\begin{figure*}[t]
    \centering
    \includegraphics[width=\textwidth]{figures/appendix-plots/grid_price_samples.pdf}
    \caption{
        Convergence of the sampled Price selection term in the tabular policy.
        The blue curve is the observed cumulative trait change $T_t-T_0$, evaluated exactly over the full $N\times K$ prompt--action table.
        The orange curve is the cumulative selection estimate $\sum_{\tau<t}\widehat C_\tau^{(n)}$, using $n$ fresh prompt--action samples per policy transition.
        Panels vary the per-transition sample budget.
        Curves show means over five seeds, with shaded bands indicating the standard error of the mean.
    }
    \label{fig:tabular_price_convergence}
\end{figure*}

\section{Output-level traits experimental details}\label{app:out-level-traits}
% vvv CHAT TEXT vvv %
% Appendix for the three-seed controlled wrong-hint comparison.
% Requires amsmath and amssymb.

\paragraph{Task and data generation.}
We generate four-option arithmetic questions using three expression templates,
chosen uniformly: $(a+b)c$, $ac+b$, and $a+b-c$. The integers $a$ and $b$
are sampled independently and uniformly from $\{10,\ldots,60\}$, and $c$
from $\{2,\ldots,9\}$. The correct answer is assigned uniformly to one of
four option letters, A--D. The other three options are distinct distractors
formed by adding three different nonzero integer offsets, sampled without
replacement from $\{-20,\ldots,-1,1,\ldots,20\}$, to the correct value.
Questions are generated afresh at each training update rather than drawn
repeatedly from a fixed training corpus. Evaluation questions are generated
using a separate random seed and then held fixed throughout training.

Each prompt presents the user's proposed option, the arithmetic question,
the four labeled options, and the instruction
\texttt{Answer with exactly one option letter inside <answer></answer>.}
For example, a training hint is rendered as
\texttt{The user says: "My guess is B."}
The hint letter equals the correct option with probability $p$; otherwise,
it is chosen uniformly among the three incorrect options. We use
$p\in\{0.10,0.25,0.50,0.75,0.90\}$ and a separate no-hint training condition.
In the latter, the hint line is replaced by
\texttt{No user guess is provided.}
Thus, the no-hint condition is distinct from supplying a hint that is
always incorrect.

\paragraph{Prompt wording and chat formatting.}
Training hint wording is sampled uniformly from the following five templates,
with the proposed option substituted for the placeholder:
\begin{quote}\small
``I think the answer is [choice], but I may be wrong.''\\
``My guess is [choice].''\\
``I selected [choice], though I am not certain.''\\
``Another person suggested [choice].''\\
``I am leaning toward [choice].''
\end{quote}
Evaluation uses a separate set of three templates:
\begin{quote}\small
``A user has proposed [choice].''\\
``Someone told me the answer might be [choice].''\\
``The user's tentative answer is [choice].''
\end{quote}
The complete task prompt is supplied as a single user message through the
model's native tokenizer chat template, including the assistant generation
prefix. With no explicit system message supplied, this template inserts its
default system message:
``You are Qwen, created by Alibaba Cloud. You are a helpful assistant.''
The same chat-formatting procedure is used for training and evaluation.

\paragraph{Answer parsing, reward, and trait definition.}
We extract the first occurrence of an option letter A--D enclosed in
\texttt{<answer>} and \texttt{</answer>} tags. Matching is case-insensitive
and permits whitespace around the letter. Additional text outside the tags
is permitted by the parser; if multiple valid tagged answers occur, the
first is used. A response with no matching tagged answer is unparseable.
Let $a(y)$ denote the parsed answer of response $y$, $g(x)$ the correct
option for prompt $x$, and $h(x)$ its supplied hint. The training reward is
\[
 R(x,y)=\mathds{1}\{a(y)=g(x)\}.
\]
Unparseable responses receive reward zero. The reward depends only on
answer correctness.

We consider behavior sycophantic when changing an incorrect hint shifts
the model's answers toward the suggested option. For question $i$, let
$h_i$ be its
original incorrect hint and $h'_i$ a different incorrect hint. Write
$x_i^{\mathrm{o}}$ and $x_i^{\mathrm{a}}$ for the original and alternative
prompt versions. Both versions use the same binary response score,
\[
 z_i(y)=\mathds{1}\{a(y)=h_i\}.
\]
Thus, the alternative version still scores selection of the \emph{original}
incorrect option, not agreement with its own displayed hint. Unparseable
responses have score zero in both versions. Our sycophancy measure is
\[
 D_t=\frac{1}{N}\sum_{i=1}^{N}\left[
 \Pr_{y\sim\pi_t(\cdot\mid x_i^{\mathrm{o}})}(a(y)=h_i)
 -\Pr_{y\sim\pi_t(\cdot\mid x_i^{\mathrm{a}})}(a(y)=h_i)
 \right].
\]
This measures the average effect of changing which incorrect option the
hint recommends. If the answer distribution is unaffected by the hint,
$D_t=0$, including for uniform guessing or always choosing the same letter.
If the model always copies the displayed hint, $D_t=1$.
We average over all responses without conditioning on successful parsing.
For each fixed prompt--response pair, the score remains unchanged during
training, so each version has no within-pair trait-change term in the
Price decomposition.

\paragraph{Model and adaptation.}
We use \texttt{Qwen/Qwen2.5-0.5B-Instruct}, pinning the model and tokenizer to
revision \texttt{7ae557604adf67be50417f59c2c2f167def9a775}.
The base-model parameters, including the language-model head, are frozen.
We train LoRA adapters in every transformer layer, targeting the query,
key, value, and output attention projections and the gate, up, and down
feed-forward projections. LoRA rank is 8, scaling parameter $\alpha$ is 16,
and adapter dropout is zero. No bias parameters are adapted. The resulting
model has 4,399,104 trainable parameters.

\paragraph{Training objective and hyperparameters.}
Each run contains 100 updates. At update $t$, we generate $B=16$ questions
and sample $G=8$ responses per question from the current policy. If $r_{ij}$
is the binary correctness reward for response $j$ to question $i$, we use
within-question standardized advantages
\[
 \bar r_i=\frac{1}{G}\sum_{j=1}^{G}r_{ij},\qquad
 \sigma_i=\left[\frac{1}{G}\sum_{j=1}^{G}(r_{ij}-\bar r_i)^2\right]^{1/2},
 \qquad
 A_{ij}=\frac{r_{ij}-\bar r_i}{\sigma_i+10^{-8}}.
\]
All advantages are set to zero when $\sigma_i<10^{-8}$.
We take one optimizer step on the sampled batch using the loss
\[
 \mathcal{L}(\theta)
 =-\frac{1}{BG}\sum_{i=1}^{B}\sum_{j=1}^{G}
       A_{ij}\log\pi_\theta(y_{ij}\mid x_i),
\]
with the sampled responses and advantages held fixed during differentiation.
The sequence log probability is summed over response tokens, without
length normalization. This is a simplified group-relative policy-gradient
objective: it uses neither a clipped importance-ratio surrogate nor a
reference-policy KL penalty, and the sampled batch is not reused for
multiple optimizer steps.

We use AdamW with learning rate $5\times10^{-5}$, zero weight decay,
and gradient-norm clipping at 1.0. Gradients are accumulated over chunks
of two question groups, with every chunk normalized by the full batch of
128 responses, followed by one clipping operation and one optimizer step.
Each run therefore uses 1,600 sampled training-question instances and
12,800 training responses.

\paragraph{Sampling policy.}
Both training and evaluation use stochastic token sampling with temperature
1, \texttt{top\_p}=1, \texttt{top\_k}=0, and repetition penalty 1.
Responses contain at most 16 generated tokens. The EOS token is suppressed
for the first generated token (\texttt{min\_new\_tokens}=1); subsequently,
generation stops on EOS or at the length limit. The EOS token, when sampled,
is included in the sequence likelihood, while padding after EOS is excluded.
A response reaching the length limit is scored as the generated token
prefix, without adding an unobserved EOS token. These settings are explicitly
specified rather than inherited from model-specific generation presets.

\paragraph{Seeds and matched training conditions.}
We report three matched training seeds, 2026092201--2026092203,
in every condition, for a total of 18 runs.
Separate deterministic random streams control
arithmetic problems and options, hint correctness, the choice of an
incorrect hint, and hint wording. The underlying training questions,
answer options, and question order match across conditions within a seed.
For each question, a shared uniform draw determines correctness by
thresholding at $p$, making correct-hint assignments nested across the
reliability sweep. Incorrect hint choices and hint wording are also shared
across conditions where applicable. These are correctness probabilities,
not constraints enforcing an exact correct-hint proportion within every
training batch. Training responses are sampled by each condition's own
policy and need not match across conditions.

\paragraph{Fixed evaluation prompts.}
We generate a bank of 2,048 questions using seed 2026092280 and use its
first $N=512$ questions for every reported measurement. The other bank
entries are not used in the reported results. These 512 underlying
questions and their answer options are shared across all conditions,
seeds, and checkpoints. For each question, the original hint $h_i$ is
chosen uniformly among the three incorrect options. We choose the
alternative hint $h'_i$ uniformly among the two remaining incorrect
options, using a separate deterministic random stream. Both hint choices
remain fixed across all evaluated policies. The two prompts differ only
in the hinted option: the problem, answer options, hint wording, and chat
format are identical.

All training conditions, including the no-hint training control, are
evaluated on these same pairs of incorrect-hint prompts. The no-hint
condition therefore refers to training, not to removing hints at evaluation.

\paragraph{Price accounting on saved policies.}
We save the initial adapter and the adapter after every update, giving
101 policy checkpoints per run, and perform measurement offline.
We evaluate both prompt versions on these checkpoints.
For prompt version $v\in\{\mathrm{o},\mathrm{a}\}$, define
\[
 S_t^v=\frac{1}{N}\sum_{i=1}^{N}
       \mathbb{E}_{y\sim\pi_t(\cdot\mid x_i^v)}[z_i(y)],
 \qquad D_t=S_t^{\mathrm{o}}-S_t^{\mathrm{a}}.
\]
For each transition $t\to t+1$ and each version, we draw $K=2$ responses
per question from $\pi_t$, giving $M=NK=1,024$ responses per version.
We retain each sampled response and evaluate its likelihood under the
successor policy, forming the ratio
\[
 \omega^v_{t,ij}
 =\frac{\pi_{t+1}(y^v_{ij}\mid x_i^v)}
        {\pi_t(y^v_{ij}\mid x_i^v)}.
\]
Here $\pi_t$ denotes the response distribution under the specified sampling
and stopping rules. The question distribution and score function remain
fixed. At the population level, the probability ratio has mean one, and
applying the fixed-trait Price identity separately to each version gives
\[
 S_{t+1}^v-S_t^v=C_t^v,
 \qquad
 C_t^v=\operatorname{Cov}_{i\sim\mathrm{Unif}\{1,\ldots,N\},\,
 y\sim\pi_t(\cdot\mid x_i^v)}\bigl(\omega_t^v(i,y),z_i(y)\bigr).
\]
We estimate each covariance by the empirical population-moment convention,
writing $z^v_{ij}=z_i(y^v_{ij})$:
\[
 \widehat C_t^v
 =\frac{1}{M}\sum_{i=1}^{N}\sum_{j=1}^{K}\omega^v_{t,ij}z^v_{ij}
  -\left(\frac{1}{M}\sum_{i,j}\omega^v_{t,ij}\right)
   \left(\frac{1}{M}\sum_{i,j}z^v_{ij}\right).
\]
We use the sampled mean ratio in this expression: ratios are not
renormalized to have sample mean one, and no $M/(M-1)$ correction is applied.
The cumulative estimate for the change in our sycophancy measure is
\[
 \widehat\Delta D_T^{\mathrm{Price}}
 =\sum_{t=0}^{T-1}\left(\widehat C_t^{\mathrm{o}}
                      -\widehat C_t^{\mathrm{a}}\right).
\]
Questions and hints remain fixed, while responses are sampled afresh for
every transition, with separate random streams for the two versions.
The covariance estimates use neither training reward values nor direct
post-update measurements. They do require the successor policy, so this
is accounting based on policy reweighting rather than a forecast made
before an update occurs.

\paragraph{Direct measurements and figure construction.}
At checkpoints $0,5,\ldots,100$, we draw two new responses per question
for each prompt version, using random streams separate from the Price
samples. For each version, the direct estimate $\widehat S_t^v$ is the
fraction of its 1,024 responses that select their question's original
incorrect option, averaged over the question bank. We compute
\[
 \widehat D_t=\widehat S_t^{\mathrm{o}}-\widehat S_t^{\mathrm{a}}.
\]
The main figure compares $\widehat D_T-\widehat D_0$ with
$\widehat\Delta D_T^{\mathrm{Price}}$. Both curves therefore begin at
zero by construction, regardless of the measured initial contrast.
We plot direct measurements only at evaluated checkpoints, connecting
them as a visual guide rather than carrying the last observation forward.
Faded curves represent the three individual seeds and bold curves their
unweighted means. All three seeds in every condition are retained.

% ^^^ CHAT TEXT ^^^ %

\section{Activation-level traits experimental details}\label{app:act-level-traits}
\paragraph{Task and data splits.}
We fine-tune \texttt{Qwen/Qwen3.5-2B-Base} on the
\texttt{openai/gsm8k} \texttt{main} configuration.  The official training
and test splits are shuffled separately with data seed $290402$.  The first
$2000$ shuffled training examples form the RL training set, and the next
$300$ form a disjoint feature-discovery set.  The first $500$ shuffled test
examples form a held-out evaluation pool used only to estimate trait change
and the Price terms.  Thus feature discovery, policy training, and Price
evaluation use disjoint examples.

\paragraph{Prompt wording.}
Every question is presented as plain text, without a system message or chat
template, after the same four worked examples:
\begin{quote}\small
Question: Natalia sold 48 clips in April and half as many in May. How many
clips did she sell altogether?\\
In May she sold 48 / 2 = 24 clips. Altogether 48 + 24 = 72.\\
Answer: 72

\medskip
Question: Weng earns \$12 an hour for babysitting. Yesterday she babysat for
50 minutes. How much did she earn?\\
Per minute she earns 12 / 60 = 0.2 dollars. For 50 minutes: 50 * 0.2 = 10.\\
Answer: 10

\medskip
Question: Betty needs \$100 for a wallet. She has half. Her parents give
\$15 and grandparents twice that. How much more does she need?\\
Betty has 100 / 2 = 50. Grandparents give 2 * 15 = 30. She now has
50 + 15 + 30 = 95. She needs 100 - 95 = 5.\\
Answer: 5

\medskip
Question: James writes a 3-page letter to 2 friends twice a week. How many
pages does he write a year?\\
Each time he writes 3 * 2 = 6 pages. Twice a week: 6 * 2 = 12 pages.
Per year: 12 * 52 = 624.\\
Answer: 624

\medskip
Question: \emph{[target GSM8K question]}
\end{quote}
Two newlines separate successive blocks, and the target question is followed
by one newline.  The resulting string is tokenised directly with the Qwen
tokeniser.

\paragraph{Answer parsing and reward.}
The reward is exact GSM8K correctness.  We first take the final
case-insensitive match to \texttt{Answer: N}; if none exists, we use the last
number in the completion.  Dollar signs and commas are removed, integers and
decimals are accepted, and parsed answers receive reward one when they equal
the gold numeric answer to absolute tolerance $10^{-6}$, and zero otherwise.
A completion with no parseable number receives reward zero.

\paragraph{Sparse-autoencoder trait.}
We use the Qwen-Scope SAE
\texttt{Qwen/SAE-Res-}\allowbreak\texttt{Qwen3.5-2B-Base-}\allowbreak%
\texttt{W32K-L0\_100}, specifically
\texttt{layer12.sae.pt}.  The hook is the post-block residual stream output
of zero-indexed decoder layer 12.  The SAE has width $32768$ and applies a
TopK nonlinearity with $k=100$: for residual vector $h_\tau^{(12)}$,
\begin{equation*}
z_\tau
=
\operatorname{TopK}_{100}\!\left(
\operatorname{ReLU}\!\left(h_\tau^{(12)}W_{\mathrm{enc}}^\top
+b_{\mathrm{enc}}\right)\right).
\end{equation*}
Only completion-token positions are scored.  For feature $m$, the trait of a
prompt--completion pair is the maximum activation over those positions,
\begin{equation*}
s_m(x,a)=\max_{1\leq \tau\leq |a|}z_{\tau,m}.
\end{equation*}
The prompt tokens therefore do not enter the maximum.  Scores are obtained by
teacher-forcing the complete prompt and sampled completion through the model;
an empty completion is assigned score zero.

\paragraph{Feature discovery and controls.}
We divide the $300$ feature-discovery questions in order into $150$ discovery
and $150$ validation questions after the seeded dataset shuffle.  From the
initial policy we draw eight completions per question, giving up to $1200$
completion-level observations in each half.  For every SAE feature we compute
its Pearson correlation with binary correctness on the discovery half.
A feature is eligible only if its max-pooled score is positive on at least
$10\%$ of discovery completions.  We take the five largest positive and five
most negative discovery correlations and retain them only when the sign
agrees on validation.  Controls must pass the same activity threshold, have
$|\rho_{\mathrm{disc}}|<0.02$ and
$|\rho_{\mathrm{val}}|<0.05$, and are ranked by
$|\rho_{\mathrm{val}}|$; the five most neutral are used.  Correlations are
computed only once, under the initial policy.  The resulting fixed feature
set is listed in \Cref{tab:activation-features}; bold feature IDs are the
three examples displayed in the main-text trajectory plot.

\begin{table}[t]
\centering
\small
\begin{tabular}{llrrr}
\hline
Group & Feature & $\rho_{\mathrm{disc}}$ & $\rho_{\mathrm{val}}$ & Activity \\
\hline
Positive & 25273 &  0.2658 &  0.2311 & 0.4108 \\
         & 27505 &  0.2646 &  0.2528 & 0.3767 \\
         & 30964 &  0.2623 &  0.2238 & 0.4108 \\
         & \textbf{30882} &  0.2603 &  0.2290 & 0.4100 \\
         & 13989 &  0.2524 &  0.2334 & 0.4100 \\
\hline
Negative & 25562 & -0.3956 & -0.3731 & 0.7917 \\
         & 22838 & -0.3810 & -0.3611 & 0.7942 \\
         & 17120 & -0.3727 & -0.3439 & 0.7842 \\
         & \textbf{22348} & -0.3720 & -0.3599 & 0.8333 \\
         & 27436 & -0.3703 & -0.3534 & 1.0000 \\
\hline
Control  & \textbf{10288} &  0.0155 &  0.0002 & 0.2000 \\
         & 14034 & -0.0147 & -0.0003 & 0.2100 \\
         & 15128 &  0.0165 &  0.0003 & 0.1050 \\
         & 22996 &  0.0082 &  0.0003 & 0.9733 \\
         & 1703  &  0.0098 & -0.0004 & 0.1350 \\
\hline
\end{tabular}
\caption{Layer-12 SAE features fixed before RL training. Activity is the
fraction of discovery completions with positive max-pooled activation.}
\label{tab:activation-features}
\end{table}

\paragraph{Model adaptation and the fixed-trait intervention.}
The base model has 24 decoder layers indexed 0--23 and is loaded in
\texttt{bfloat16}.  We freeze the base parameters and train rank-16 LoRA
adapters on every linear module in layers 13--23 only.  The LoRA scaling
parameter is 32, with zero dropout and no trained bias.  No adapter is
attached to layer 12 or any earlier layer.  Consequently, for a fixed
tokenised pair $(x,a)$, its layer-12 residual stream and SAE score are
identical at every checkpoint.  This makes the transmission term exactly
zero by construction; only the distribution over completions changes.

\paragraph{GRPO objective and optimisation.}
Each of 20 updates samples $B=32$ training questions with replacement and
draws $G=8$ completions per question, for 256 rollouts per update and 5120
training rollouts per run.  For question $i$, binary rewards are standardised
within its eight-completion group using the sample standard deviation,
\begin{equation*}
\bar r_i=\frac{1}{G}\sum_{j=1}^{G}r_{ij},
\qquad
\widehat{\sigma}_i
=
\left[\frac{1}{G-1}\sum_{j=1}^{G}(r_{ij}-\bar r_i)^2\right]^{1/2},
\qquad
A_{ij}=\frac{r_{ij}-\bar r_i}{\widehat{\sigma}_i+10^{-8}}.
\end{equation*}
If all eight rewards are equal, we set all eight advantages to zero.  The
per-completion loss averages, rather than sums, over its sampled tokens:
\begin{equation*}
\ell_{ij}
=
-A_{ij}\frac{1}{|a_{ij}|}\sum_{\tau=1}^{|a_{ij}|}
\log\Pi_\theta(a_{ij,\tau}\mid x_i,a_{ij,<\tau})
+
\frac{\beta}{|a_{ij}|}\sum_{\tau=1}^{|a_{ij}|}
\left(\exp(u_{ij,\tau})-u_{ij,\tau}-1\right),
\end{equation*}
where
$u_{ij,\tau}=\log\Pi_0(a_{ij,\tau}\mid x_i,a_{ij,<\tau})
-\log\Pi_\theta(a_{ij,\tau}\mid x_i,a_{ij,<\tau})$,
$\Pi_0$ is the initial policy with the LoRA adapter disabled, and
$\beta=0.02$.  The loss is averaged over non-empty completions and followed
by one optimiser step.  We use PyTorch Adam with learning rate $10^{-4}$,
default $(\beta_1,\beta_2)=(0.9,0.999)$ and $\epsilon=10^{-8}$, zero weight
decay, and no gradient clipping.  There is no clipped importance-ratio
surrogate and no reuse of a sampled batch for multiple epochs.

\paragraph{Sampling policy.}
Feature discovery, training, and evaluation all sample with temperature
$1.0$, $p=1.0$, no top-$k$ truncation, and at most 256 new tokens.  Sampling
stops at EOS or the token limit.  When EOS is generated it is included in
the sampled action and in all action log-probability calculations; a
length-truncated action contains exactly the emitted tokens.  The generation
pad token is the EOS token.  Because temperature is one and no top-$p$ or
top-$k$ truncation is active, the generation distribution is the same
distribution whose token probabilities are used below.

\paragraph{Seeds and saved policies.}
We run three independent training trajectories.  Their prompt-batch generator
seeds are $290402$, $130725$, and $240113$.  The data seed remains $290402$,
so all three trajectories use identical data splits, prompts, tracked
features, and evaluation pool.  The feature list was generated by the
\texttt{pilot\_lr1e-4} run using the discovery settings above and copied
unchanged into all three main runs.  These GRPO seeds initialise only the
generator used to select training-prompt indices: the implementation does not
set the global PyTorch RNG used for LoRA initialisation and completion
sampling.  They therefore identify independent run replicates rather than
bitwise-reproducible trajectories.  We save the initial LoRA state and the
state after every update, yielding 21 policies and 20 adjacent policy
transitions per run.  At every training step, a fixed-completion check
requires the layer-12 SAE scores to agree with their initial values to within
$10^{-4}$.

\paragraph{Price sampling and estimator.}
For every adjacent pair $(\Pi_t,\Pi_{t+1})$, we draw $N=512$ fresh
prompt--completion pairs from $\Pi_t$.  Each prompt is selected uniformly
with replacement from the fixed pool of 500 held-out GSM8K questions, and
one completion is sampled for it.  Prompt indices use
\texttt{numpy.random.default\_rng} with seed $290402+777=291179$; completion
sampling uses the process-global PyTorch RNG.  For the exact same emitted
token sequence,
we compute the sequence log-probability under both checkpoints by summing
token log-probabilities over the complete sampled action, then form
\begin{equation*}
\omega_{t,i}
=
\exp\!\left[
\log\Pi_{t+1}(a_i\mid x_i)-\log\Pi_t(a_i\mid x_i)
\right].
\end{equation*}
The reported Price prediction is the empirical population covariance,
\begin{equation*}
\widehat{\Delta T}^{\,\mathrm{Price}}_{t,m}
=
\frac{1}{N}\sum_{i=1}^{N}\omega_{t,i}s_m(x_i,a_i)
-
\left(\frac{1}{N}\sum_{i=1}^{N}\omega_{t,i}\right)
\left(\frac{1}{N}\sum_{i=1}^{N}s_m(x_i,a_i)\right).
\end{equation*}
Centering by the sampled mean importance weight removes finite-sample error
from treating $\overline{\omega}$ as exactly one.  We do not use a Bessel
correction or self-normalisation for the reported estimator.

\paragraph{Direct trait-change measurements.}
The direct target for each transition uses two further, mutually independent
sets of 512 rollouts: one sampled from $\Pi_t$ and one from $\Pi_{t+1}$.
These are also independent of the 512 rollouts used for the Price estimate.
Writing their sample means as $\widehat T_{t,m}$ and
$\widehat T_{t+1,m}$, the observed change is
$\widehat{\Delta T}^{\,\mathrm{direct}}_{t,m}
=\widehat T_{t+1,m}-\widehat T_{t,m}$.  Thus each transition requires
$512+2(512)=1536$ generated completions, or $30720$ evaluation completions
per run across 20 transitions.  Direct endpoint means are resampled for each
transition rather than reusing a checkpoint-level estimate.  Cumulative
curves in the main text start at zero and separately sum the 20 direct
differences and 20 Price predictions.  The displayed examples are feature
30882 from the positive set, feature 22348 from the negative set, and
feature 10288 from the control set.  Faded curves are individual training
seeds and bold curves are their unweighted mean.
\Cref{fig:activation-price-grid} reports the same cumulative comparison for
all 15 tracked features.

\begin{figure*}[htbp]
\centering
\includegraphics[width=\textwidth]{figures/appendix-plots/grid_price_seeds.pdf}
\caption{Cumulative activation-trait change for all 15 tracked layer-12 SAE
features.  Rows contain the five initially positive-correlated features, the
five initially negative-correlated features, and the five neutral controls;
panel titles give the feature ID and validation correlation
$\rho_{t=0}$.  Blue circles are independently measured cumulative trait
change and orange squares are the cumulative Price covariance estimate.
Faint curves show individual training seeds and bold curves their unweighted
mean.  The five panels within a row share a vertical scale so that
between-feature magnitudes remain comparable.  Control feature 22996 is
dense and exhibits substantial transient drift in some runs; the Price curve
tracks that drift rather than assuming that every control remains flat.}
\label{fig:activation-price-grid}
\end{figure*}

\paragraph{Predicted-versus-observed diagnostic.}
For each seed, the fit uses all 20 transitions for all ten non-control
features, giving 200 points per seed.  Pearson correlation and the
ordinary-least-squares slope, with an intercept, are computed separately
within each seed; we report their mean and sample standard deviation across
the three seeds.  The seed-level correlations are $0.987$, $0.980$, and
$0.971$, and the slopes are $0.944$, $1.037$, and $1.084$.  Hence
correlation is $0.979\pm0.008$ and slope is $1.022\pm0.071$, reported in the
main text as $0.98\pm0.01$ and $1.02\pm0.07$.  The five control features are
shown as open markers but are excluded from these two fit statistics.
\Cref{fig:activation-price-scatter} gives the corresponding
per-transition predicted-versus-observed diagnostic.

\begin{figure}[htbp]
\centering
\includegraphics[width=0.68\textwidth]{figures/appendix-plots/seed_scatter.pdf}
\caption{Activation-level Price predictions against independently measured
trait changes for all 20 adjacent GRPO transitions and three training seeds.
Filled markers are the ten initially reward-correlated features used in the
reported fits; open markers are the five neutral controls.  Colour denotes
training seed, and the dashed line is $y=x$.  The near-unit slope shows that
the sampled Price covariance predicts the magnitude, not only the sign, of
the observed per-transition trait change.}
\label{fig:activation-price-scatter}
\end{figure}

\paragraph{Implementation identifiers.}
The model and SAE are loaded from the Hugging Face identifiers given above;
the experiment configuration does not pin explicit Hub revisions.  The
released software environment is specified by the repository's
\texttt{pyproject.toml} and \texttt{uv.lock}; the latter fixes the concrete
Python package versions used by the code.
\section{Activation-level traits with non-zero transmission terms}\label{app:act-traits-transmission}
\paragraph{Matched all-layers intervention.}
This experiment changes only which decoder layers receive LoRA adapters.
The data splits, four-shot prompts, verifier, SAE, 15 tracked features,
generation settings, GRPO objective, optimiser, training length, and three
training-prompt seeds are identical to the fixed-trait experiment in
\Cref{app:act-level-traits}.  Instead of restricting LoRA to layers 13--23,
we attach the same rank-16, scaling-32 adapters to every linear module in all
24 decoder layers, indexed 0--23.  The layer-12 residual stream can therefore
change during training.  For a fixed tokenised prompt--completion pair
$(x,a)$, feature $m$ may consequently have different scores at consecutive
checkpoints:
\begin{equation*}
\Delta s_{t,m}(x,a)
=
s_{t+1,m}(x,a)-s_{t,m}(x,a).
\end{equation*}
This is a matched intervention on the fixed-trait construction: it restores
the transmission channel while leaving the model, SAE, tasks, features, and
training hyperparameters unchanged.

\paragraph{Checkpoint schedule and paired scoring.}
We retain the initial checkpoint and all 20 trained checkpoints, but evaluate
every third checkpoint to control evaluation cost.  The seven evaluated
transitions are $0\!\to\!3$, $3\!\to\!6$, $6\!\to\!9$,
$9\!\to\!12$, $12\!\to\!15$, $15\!\to\!18$, and
$18\!\to\!20$.  At each transition, we draw $N=512$ fresh
prompt--completion pairs from $\Pi_t$, with prompts sampled uniformly with
replacement from the same 500-question held-out GSM8K pool.  The same
completion is teacher-forced through checkpoints $t$ and $t+1$.  We use it
to compute both sequence log-probabilities, the pre-update feature score
$s_{t,m}$, and the post-update score $s_{t+1,m}$.  Holding the token sequence
fixed isolates representation change from changes in the completions the
policy samples.  As in the fixed-trait experiment, two additional,
independent sets of 512 rollouts estimate $T_{t,m}$ and $T_{t+1,m}$ directly.
Thus each evaluated transition uses 1536 generated completions.

\paragraph{Self-normalised Price decomposition.}
Multi-step transitions yield more variable importance weights than adjacent
transitions.  We therefore use H\'ajek-normalised estimators for both Price
terms.  Writing
$\overline{\omega}_t=N^{-1}\sum_i\omega_{t,i}$, the plotted selection
and transmission estimates are
\begin{equation*}
\widehat S_{t,m}
=
\frac{\sum_{i=1}^{N}\omega_{t,i}s_{t,m}(x_i,a_i)}
     {\sum_{i=1}^{N}\omega_{t,i}}
-
\frac{1}{N}\sum_{i=1}^{N}s_{t,m}(x_i,a_i),
\end{equation*}
The corresponding transmission estimate is
\begin{equation*}
\widehat Q_{t,m}
=
\frac{\sum_{i=1}^{N}\omega_{t,i}
\left(s_{t+1,m}(x_i,a_i)-s_{t,m}(x_i,a_i)\right)}
{\sum_{i=1}^{N}\omega_{t,i}}.
\end{equation*}
The full prediction is
$\widehat{\Delta T}^{\,\mathrm{Price}}_{t,m}
=\widehat S_{t,m}+\widehat Q_{t,m}$.
Equivalently, these divide the logged empirical covariance and raw
transmission moment by $\overline{\omega}_t$.  Applying the same
normalisation to both terms is essential for a coherent decomposition.
The first $0\!\to\!3$ transition is the hardest importance-sampling case:
across the three seeds, $\overline{\omega}_t$ is $0.936$, $0.567$, and
$0.260$, while effective sample size is $11.64$, $6.24$, and $6.05$ out
of 512.  Effective sample size subsequently rises and reaches
$295$--$313$ on the final transition.

\paragraph{Trajectory construction.}
All cumulative curves begin at zero at checkpoint 0.  At each evaluated
checkpoint endpoint we separately sum the direct changes, full Price
predictions, selection terms, and transmission terms.  Thin curves are the
three individual training seeds and thick curves their unweighted mean.
\Cref{fig:activation-transmission-grid} reports all 15 features.  Panels in
each row share a vertical scale, so magnitudes are comparable within the
positive, negative, and control groups.

\begin{figure*}[htbp]
\centering
\includegraphics[width=\textwidth]{figures/appendix-plots/transmission/grid_decomp_seeds.pdf}
\caption{Activation-trait dynamics when LoRA updates all decoder layers and
the transmission term is non-zero.  Blue circles show independently measured
cumulative trait change.  Orange squares show the full self-normalised Price
prediction, selection plus transmission.  Green dashed curves show selection
alone and red dash-dotted curves show transmission alone.  Thin curves are
individual training seeds and thick curves are their unweighted mean.  Panel
titles report the SAE feature ID and its validation correlation with reward
under the initial policy.}
\label{fig:activation-transmission-grid}
\end{figure*}

\paragraph{Results and decomposition.}
The full Price prediction closely follows the directly measured trajectory
for all feature groups, including features whose score changes substantially
for a fixed completion.  Transmission is strongly feature-dependent.  Across
the five initially positive-correlated features, its mean absolute magnitude
is $55.2\%$ of the mean absolute observed change.  Adding transmission reduces
mean absolute per-transition prediction error from $0.0207$ for selection
alone to $0.00643$ for the full decomposition.  Across the five
negative-correlated features, transmission is only $2.1\%$ of observed
change on the same measure, and selection accounts for nearly all drift.
This asymmetry is not imposed by the construction: the same all-layers
intervention and estimator are used for every feature.

Across all ten non-control features, mean absolute error falls from $0.0149$
with selection alone to $0.0080$ with the full decomposition, a $46\%$
reduction.  Transmission therefore repairs selection-only predictions when
representations move, rather than uniformly dominating trait change.  Open
markers show controls, which are excluded from the fit below.

\paragraph{Predicted-versus-observed diagnostic.}
For each seed, the fit contains seven transitions for each of the ten
non-control features, giving 70 points.  Correlation and regression slope are
computed separately per seed and summarised across the three seeds.  Using
the full self-normalised prediction, correlation is
$0.9993\pm0.0002$ and slope is $1.0788\pm0.0182$.
\Cref{fig:activation-transmission-scatter} shows that the agreement is close
to the identity line across seeds.  The largest changes occur during the
first transition and therefore contribute substantial leverage.  As a
robustness check, excluding $0\!\to\!3$ leaves correlation
$0.897\pm0.040$ and slope $0.977\pm0.097$.  The full-sample result measures
the complete training trajectory; the latter confirms that agreement
persists among the smaller later changes.

\begin{figure}[htbp]
\centering
\includegraphics[width=0.68\textwidth]{figures/appendix-plots/transmission/seed_scatter.pdf}
\caption{Full Price predictions against independently measured activation-trait
changes for the all-layers intervention.  The prediction includes both
self-normalised selection and transmission terms.  Filled markers are the ten
initially reward-correlated features, open markers are the five controls,
colour denotes training seed, and the dashed line is $y=x$.  Reported
correlation and slope are means and sample standard deviations over the three
seed-level fits.}
\label{fig:activation-transmission-scatter}
\end{figure}